\section{CMP decision algorithms: proved and required forms}
\label{sec:algorithm}

Two procedures must be distinguished.  The first is an exact algorithm for an
evaluated linearized system.  The second is the desired but currently
unconstructed exact feasibility decider.

\subsection{Differential CMP}

Given numerical Jacobian blocks at a specified point \(p\), differential CMP
performs the following operations.

\begin{enumerate}
  \item For each sample, form \(C_i(p)\) and \(D_i(p)\).
  \item Compute a basis \(P_i^\perp\) for
  \(\kernel D_i(p)^{\mathsf T}\).
  \item Form the local parameter message
  \begin{equation}
    M_i=P_i^\perp C_i(p).
  \end{equation}
  \item Stack the messages and remove linearly dependent rows.
  \item Accept a proposed parameter direction \(\delta\theta\) if and only if
  the root matrix annihilates \(\delta\theta\).
\end{enumerate}

\begin{theorem}[Correctness of differential CMP]
\label{thm:differential-cmp-correct}
At a specified real point \(p\), the root test accepts
\(\delta\theta\) if and only if there exist local directions
\(\delta S^{(i)}\) satisfying all linearized equations
\begin{equation}
  C_i(p)\delta\theta+D_i(p)\delta S^{(i)}=0.
\end{equation}
\end{theorem}

\begin{proof}
Apply \cref{prop:local-linear-extension} independently to every sample and
conjoin the resulting conditions.  Equivalently, recursively apply
\cref{thm:linear-separator-factorization} along the star decomposition.
\end{proof}

This procedure assumes the point \(p\) is already supplied.  It neither tests
whether \(p\) exists nor decides the original instance.

\subsection{The required exact procedure}

An exact decision algorithm would need a finite message representation
\(R_v\) for every subtree with a denotation
\(\denote{R_v}\subseteq\R^{S_v}\).  Its operations must satisfy:

\begin{description}[leftmargin=3.4cm,style=nextline]
  \item[Leaf construction]
  \(\denote{R_v}\) is exactly the projection of the leaf's ReLU
  constraints to its separator.

  \item[Merge]
  Combining child representations denotes the intersection of their lifted
  feasible sets with the current bag constraints.

  \item[Contraction]
  Eliminating internal variables denotes existential projection, with both
  soundness and completeness.

  \item[Root test]
  The algorithm returns \textsc{yes} exactly when the root denotation is
  nonempty.

  \item[Complexity]
  Representation length and the time for every operation are polynomial in
  the original binary input length.
\end{description}

If arbitrary quantifier-free formulas over the reals are used as
representations, the first four requirements follow from standard logical
semantics.  The fifth does not follow.  If only the cotangent subspaces of
\cref{def:intrinsic-message} are used, the fifth is inexpensive but the first
and third fail by \cref{thm:cotangent-only-false}.

\begin{problem}[Polynomial semantic message]
\label{prob:polynomial-semantic-message}
Construct a representation and algorithms satisfying all five requirements
for every finite instance of \(\ExactNNTP\), or prove that no such
representation exists under an explicitly stated computational model.
\end{problem}

\subsection{No hidden root oracle}

The root test must be included in the complexity analysis.  A procedure that
propagates compact differential messages and then invokes an unrestricted
semialgebraic feasibility solver at the root has not placed the language in
\(P\); it has moved the original difficulty into the final operation.  The
same applies if a ``message entry'' is an unevaluated existential formula whose
nonemptiness remains ETR-hard.
